% Préambule commun au cours (report) et aux fiches de colle (article).
% Chargé via \IfFileExists pour fonctionner depuis la racine comme depuis colles/.

\usepackage[T1]{fontenc}
\usepackage[french]{babel}
\usepackage{amsmath, amssymb, amsthm}
\usepackage{stmaryrd}
\usepackage{geometry}
\usepackage[hidelinks]{hyperref}
\geometry{margin=2.5cm}

% --- Environnements de théorèmes -------------------------------------------
% Tous les énoncés partagent un même compteur, numéroté par chapitre lorsque
% la classe en fournit un (report), globalement sinon (article).
\theoremstyle{plain}
\ifdefined\chapter
  \newtheorem{prop}{Proposition}[chapter]
\else
  \newtheorem{prop}{Proposition}
\fi
\newtheorem{theo}[prop]{Théorème}
\newtheorem{lemme}[prop]{Lemme}
\newtheorem{coro}[prop]{Corollaire}

\theoremstyle{definition}
\newtheorem{defi}[prop]{Définition}
\newtheorem{exemple}[prop]{Exemple}

\theoremstyle{remark}
\newtheorem*{remarque}{Remarque}

% --- Macros ----------------------------------------------------------------
\newcommand{\inter}[2]{\llbracket #1,#2 \rrbracket}   % intervalle d'entiers
\newcommand{\Card}{\operatorname{Card}}
\newcommand{\Ker}{\operatorname{Ker}}
\newcommand{\Img}{\operatorname{Im}}
\newcommand{\Vect}{\operatorname{Vect}}
\newcommand{\Spec}{\operatorname{Spec}}
\newcommand{\rg}{\operatorname{rg}}
\newcommand{\Cl}{\operatorname{Cl}}
\newcommand{\Epi}{\operatorname{Epi}}
\newcommand{\id}{\operatorname{id}}
